% =============================================================================
% Seção 2 — Estado da arte
% =============================================================================
\section{Estado da arte}
\label{sec:estado_da_arte}

\subsection{Otimização aplicada ao projeto de fundações superficiais}
\label{sec:soa_fundacoes}

A automação do projeto de fundações por meio de métodos de otimização tem sido objeto de investigação sistemática nas últimas décadas. \citeonline{wang2008economic} formalizaram o dimensionamento econômico de fundações superficiais como um problema de minimização de custos, considerando explicitamente custos de escavação, concreto, reforço e execução. Em estudos posteriores, \citeonline{khajehzadeh2012gsa} aplicaram o \textit{Gravitational Search Algorithm} ao projeto econômico de fundações rasas com restrições geotécnicas e estruturais, enquanto \citeonline{gandomi2018construction} compararam técnicas modernas de inteligência de enxames na minimização do custo de fundações rasas em concreto armado. \citeonline{nigdeli2018metaheuristic}, por sua vez, compararam DE, PSO, HS, FPA e TLBO em sapatas de concreto armado, incluindo dimensões, excentricidade relativa do pilar e detalhamento de armaduras sob verificações geotécnicas e estruturais. Em conjunto, esses trabalhos indicam que metaheurísticas populacionais podem produzir soluções economicamente competitivas quando as restrições de projeto são formuladas de modo explícito.

Mais recentemente, \citeonline{kashani2020optimum} reportaram um estudo abrangente sobre o dimensionamento ótimo de fundações superficiais utilizando algoritmos evolutivos, incluindo análise de sensibilidade em relação a variáveis geotécnicas e estruturais. Em paralelo, \citeonline{waheed2022practical} e \citeonline{waheed2025optimization} propuseram ferramentas práticas de otimização para sapatas isoladas em concreto armado, com foco no equilíbrio entre rigor normativo, economia e usabilidade por profissionais. Tais trabalhos confirmam a maturidade do uso de metaheurísticas para o problema, mas deixam espaço para investigar arquiteturas assistidas por modelos substitutos, especialmente quando se deseja estudar sistematicamente o efeito do tratamento das restrições e preparar extensões com avaliações de projeto potencialmente mais custosas.

\subsection{Métodos baseados em modelos substitutos}
\label{sec:soa_surrogates}

A combinação de metaheurísticas com modelos substitutos (\textit{surrogate-assisted optimization}) constitui uma possibilidade metodológica para problemas de engenharia em que se deseja reduzir avaliações diretas, explorar incerteza preditiva ou preparar a formulação para modelos de maior custo computacional. O método \textit{Efficient Global Optimization} (EGO), introduzido por \citeonline{jones1998efficient}, popularizou o uso de modelos do tipo \textit{kriging}/processos gaussianos como aproximadores de baixo custo para funções objetivo cuja avaliação tem custo computacional elevado. A ideia fundamental é construir, a partir de um conjunto de amostras iniciais, uma superfície probabilística que estime tanto o valor esperado da função em pontos não amostrados quanto a sua incerteza. Sobre essa superfície define-se uma \textit{função de aquisição} -- tipicamente a \textit{Expected Improvement} (EI) -- que escolhe o próximo ponto a ser avaliado equilibrando explotação (regiões de baixo valor predito) e exploração (regiões de alta incerteza).

A consolidação dos métodos bayesianos de otimização foi posteriormente sistematizada por \citeonline{snoek2012practical} no contexto da seleção de hiperparâmetros de algoritmos de aprendizado de máquina, e revisada de forma abrangente por \citeonline{shahriari2016review}, que delinearam o ciclo completo do EGO genérico, suas variantes e os principais desafios práticos -- escolha do \textit{kernel}, tratamento de restrições, paralelização e amostragem em altas dimensões. \citeonline{schulz2018tutorial} oferecem um tutorial didático sobre Regressão por Processos Gaussianos (GPR), destacando a interpretação probabilística do modelo, a estrutura dos \textit{kernels} e o impacto da escolha das funções de covariância sobre o comportamento do modelo substituto. A formalização clássica do GPR é dada em \citeonline{williams1995gaussian}. Em projeto estrutural, \citeonline{mathern2021multiobjective} demonstraram uma formulação de otimização bayesiana multiobjetivo com restrições para vigas de concreto armado, explorando explicitamente um padrão comum em engenharia: objetivos de baixo custo de cálculo e restrições normativas mais custosas.

A aplicação de métodos baseados em modelos substitutos a problemas geotécnicos é recente, porém promissora. \citeonline{ahmad2021gpr} empregaram a regressão por processos gaussianos para estimar a capacidade de carga última de fundações superficiais em solos não coesivos, evidenciando que o GPR pode capturar relações não lineares entre parâmetros geotécnicos e resposta mecânica. Em linha semelhante de modelagem por dados, \citeonline{khajehzadeh2022effective} combinaram redes neurais e metaheurística para estimar capacidade de carga e otimizar sapatas, enquanto \citeonline{fattahi2025settlement} utilizaram HS e TLBO para predizer recalques a partir de variáveis como largura, pressão aplicada, SPT, embutimento relativo e razão geométrica. Esses estudos não validam diretamente a correlação empírica adotada neste artigo, mas reforçam que capacidade de carga, recalque e variabilidade geotécnica exigem calibração específica quando deixam de ser tratados como hipóteses preliminares. Em uma direção mais ampla, \citeonline{deng2026metamaterial} combinaram GPR, autoencoder e aprendizado ativo no projeto de metamateriais, ilustrando o uso de incerteza preditiva em estratégias modernas de exploração de espaços de projeto. Embora esse último trabalho não trate de fundações, ele é usado aqui apenas como inspiração metodológica para arquiteturas com modelos substitutos e aprendizado ativo em problemas de projeto computacional.

\subsection{Tratamento de restrições e comparação entre metaheurísticas}
\label{sec:soa_restricoes_meta}

O tratamento de restrições é um aspecto central da otimização aplicada à engenharia. Estratégias baseadas em penalização exterior são simples de implementar, mas dependem de uma escolha cuidadosa do fator de penalidade: valores reduzidos podem favorecer soluções inviáveis, enquanto valores excessivos ampliam a escala da função pseudo-objetivo e podem dificultar a modelagem estatística da resposta. No contexto de modelos substitutos, esse efeito é particularmente relevante, pois o GPR precisa aproximar simultaneamente a variação física associada ao volume e a variação artificial introduzida pelas penalizações.

A literatura sobre comparação de metaheurísticas é igualmente extensa. \citeonline{morales2020balance} discutem o equilíbrio entre exploração e intensificação em diferentes algoritmos populacionais, demonstrando que tal equilíbrio é decisivo para o desempenho prático. \citeonline{abualigah2021arithmetic} apresentam o \textit{Arithmetic Optimization Algorithm}, exemplificando uma das muitas propostas recentes de novos operadores. \citeonline{gomes2018probabilistic} propõem uma métrica probabilística para comparação rigorosa entre metaheurísticas em problemas de engenharia, ressaltando a importância de avaliar não apenas o melhor valor obtido, mas também a robustez do algoritmo ao longo de múltiplas execuções.

\subsection{Lacuna identificada}
\label{sec:soa_lacuna}

Da revisão acima, identificam-se dois conjuntos de lacunas que motivam a abordagem adotada nesta pesquisa:
\begin{itemize}
    \item \textbf{Integração metodológica.} Embora existam tanto trabalhos sobre dimensionamento ótimo de fundações por metaheurísticas quanto trabalhos sobre otimização global assistida por modelos substitutos em outros domínios da engenharia, ainda precisa ser verificado em revisão bibliográfica ampliada se há estudos que unificam essas duas frentes em um único ambiente computacional voltado especificamente ao projeto de sapatas isoladas compatível com o contexto normativo brasileiro (NBR~6118 e NBR~6122).
    \item \textbf{Posicionamento como variável.} Nos trabalhos revisados sobre otimização de sapatas isoladas, o foco principal recai sobre dimensões, custo, armadura ou geometria da fundação, enquanto a posição dos pilares/fundações tende a ser tratada como dado de entrada. A consideração explícita do problema de \textit{empacotamento} -- com ausência de sobreposição, fronteira do lote, proximidade construtiva e eventual decisão por sapatas associadas -- permanece como lacuna para uma etapa posterior desta pesquisa.
\end{itemize}

A frente abordada neste artigo concentra-se na primeira lacuna. A incorporação explícita do problema de empacotamento é apresentada como trabalho futuro na Seção~\ref{sec:conclusoes}.
